\documentclass[../../../include/open-logic-section]{subfiles}

\begin{document}

\olfileid[tr]{sth}{story}{approach}
\olsection{Birikimli-Yinelemeli Yaklaşım}

Kümeleri evrelere nasıl ayıracağımızı biraz daha ayrıntılı şöyle ifade
edebiliriz:
\begin{quote}
  Kümeler \emph{evreler} hâlinde oluşturulur. Her $S$ evresi için $S$'den
  \emph{önce} gelen belirli evreler vardır. $S$ evresinde, $S$'den önceki
  evrelerde oluşturulmuş kümelerden meydana gelen her koleksiyon bir küme
  hâline getirilir. Evrelerde oluşturulan kümeler dışında hiçbir küme yoktur.
  \citep[p.~323]{Shoenfield:AST}
\end{quote}
Bu, \emph{birikimli-yinelemeli küme anlayışının} bir taslağıdır.
\olref[sth][][]{part} içinde sunduğumuz biçimsel küme teorisinin temelini
oluşturacaktır.

Bunu biraz daha ayrıntılı inceleyelim. Shoenfield'ın süreci betimleyişine
göre her evrede, önceki evrelerden elimizde bulunan kümelerden yeni kümeler
oluştururuz. Dolayısıyla başlangıç evresinde, yani $0$ evresinde daha
\emph{önceki} hiçbir evre ve bu yüzden önceki evrelerden elimizde bulunan
hiçbir küme yoktur.\footnote{Neden bir ilk evrenin \emph{var olduğunu}
varsayalım? \olref[z][story]{sec} içindeki $\stagesord{}$ dipnotuna bakın.}
Bu nedenle yalnızca bir küme oluştururuz: hiç !!{element}s{}ı olmayan
$\emptyset$. $1$ evresinde önceki evrelerden yalnızca bir küme elimizdedir;
bu nedenle tek yeni küme $\{\emptyset\}$'tir. $2$ evresinde önceki evrelerden
iki küme elimizdedir ve iki yeni küme, $\{\{\emptyset\}\}$ ile
$\{\emptyset,\{\emptyset\}\}$'i oluştururuz. $3$ evresinde önceki
evrelerden dört küme vardır; dolayısıyla on iki yeni küme oluştururuz
\ldots{} Böylece kümelerin birikimli-yinelemeli resmi, sayılar evreleri
göstermek üzere, biraz şöyle görünür:
\begin{center}
  \begin{tikzpicture}[scale=0.6]
  \tikzset{cut_here/.style={densely dotted}}
  \draw (-4,6) -- (-1, 0)-- (2,6);
  \draw[cut_here] (-5,8)--(-4,6);
  \draw[cut_here] (2,6)--(3,8);
  \draw(-1.5, 1)--(-0.5, 1);
  \draw(-2, 2)--(0, 2);
  \draw(-2.5, 3)--(0.5,3);
  \draw(-3, 4)--(1, 4);
  \draw(-3.5, 5)--(1.5, 5);
  \draw(-4, 6)--(2, 6);
  \node[label] at (0, 0) {\small 0};
  \node[label] at (0.5, 1) {\small 1};
  \node[label] at (1, 2) {\small 2};
  \node[label] at (1.5, 3) {\small 3};
  \node[label] at (2, 4) {\small 4};
  \node[label] at (2.5, 5) {\small 5};
  \node[label] at (3, 6) {\small 6};
  \end{tikzpicture}
\end{center}
Peki bu anlatıyı neden benimsemeliyiz?

Bir neden, bunun hoş ve ele alınabilir bir anlatı olmasıdır. En açık
anlatının, yani Naif Küme Oluşturma'nın çöküşünden sonra hoş bir şeye
ihtiyacımız vardır.

Fakat anlatı \emph{yalnızca} hoş değildir. Buna dayalı her küme teorisinin
\emph{tutarlı} olacağına inanmak için iyi bir nedenimiz vardır. Şöyle:

Birikimli-yinelemeli küme anlayışında kümeleri evrelerde oluştururuz ve
!!{element}s{}ı daha \emph{önceden} mevcut nesneler olmalıdır. Dolayısıyla her
$S$ evresi için şu kümeyi oluşturabiliriz:
\[
  R_S=\Setabs{x}{x\notin x\text{ ve $x$, $S$'den önce mevcuttu}}
\]
Russell Paradoksu'nun ispatındaki akıl yürütme şimdi $R_S$'nin kendisinin
$S$ evresinden önce mevcut olmadığını gösterir. Bu bir çelişki değildir.
Üstelik birikimli-yinelemeli küme anlayışını benimsiyorsak Russell kümesinin
kendisini oluşturabilmeyi zaten \emph{beklememeliyiz}. Çünkü bu, ``ileride
herhangi bir zamanda mevcut olacak'' ve kendisinin elemanı olmayan bütün
kümelerin kümesi olurdu. Kısacası Russell kümesini (kanıtlanabilir biçimde)
oluşturamamamız birikimli-yinelemeli anlatı açısından \emph{şaşırtıcı}
değildir; tam da \emph{öngöreceğimiz} şeydir.

%In one sense, then, the cumulative-iterative conception of set yields a response to Russell's Paradox which is rather like the \emph{predicativist}'s response. After all, the predicativist said that the collection of all non-self-membered sets$_n$ is a set$_{n+1}$, and this set$_n$ will not be self-membered. (Indeed, most predicativists will treat it as \emph{ungrammatical} to try to ask whether a set is self-membered.)
%
%But there is an important difference between the cumulative-iterative approach and the predicativist approach: the cumulative-iterative approach treats all of our entities as being \emph{of the same kind}. The predicativist will presumably have an empty set$_0$ (i.e., a set$_0$ with no !!{element}s), and an empty set$_1$ (i.e., a set$_1$ with no !!{element}s), and these will be \emph{different entities}. But on the cumulative-iterative approach, there is just one empty set, $\emptyset$, and it will be available at every stage of the hierarchy.

%Indeed, the Axiom of Extensionality, stated at the start of this chapter, will hold true of the elements in this hierarchy (so that there can be \emph{only one} ``empty'' set). But I do not intend to use the cumulative-iterative conception to justify Extensionality (see \cite{Boolos1971}). Again: I suggest that we should accept Extensionality, just because we are interested in extensional collections.

\end{document}
